\documentclass[11pt]{article}
\newcommand{\DocShort}{R\'ubrica}
\usepackage{fontspec}
\usepackage[spanish,es-noshorthands]{babel}
\decimalpoint
\usepackage{geometry}
\geometry{letterpaper, top=2.5cm, bottom=2.5cm, left=2.8cm, right=2.8cm}
\usepackage[expansion=false]{microtype}
\usepackage{parskip}
\setlength{\parskip}{6pt}
\usepackage{amsmath,amssymb}
\usepackage{booktabs}
\usepackage{array}
\usepackage{longtable}
\usepackage{caption}
\usepackage[dvipsnames,table]{xcolor}
\usepackage{enumitem}
\setlist[itemize]{topsep=2pt, itemsep=2pt, parsep=0pt}
\setlist[enumerate]{topsep=2pt, itemsep=2pt, parsep=0pt}
\usepackage{fancyhdr}
\usepackage[hidelinks,pdfencoding=auto]{hyperref}
\usepackage{titlesec}
\usepackage{tcolorbox}
\tcbuselibrary{breakable}
\usepackage{pdfpages}

\definecolor{darkblue}{HTML}{1B3A5C}
\definecolor{medblue}{HTML}{2E6DA4}
\definecolor{lightblue}{HTML}{D6E8F7}
\definecolor{teal}{HTML}{1A7A6E}
\definecolor{lightteal}{HTML}{D0EDEA}
\definecolor{lightgrey}{HTML}{F5F5F5}
\definecolor{midgrey}{HTML}{6F6F6F}
\definecolor{warnyellow}{HTML}{FFF3CD}

\titleformat{\section}
{\large\bfseries\color{darkblue}}
{\thesection}{0.75em}{}[{\color{medblue}\titlerule[0.8pt]}]
\titleformat{\subsection}
{\normalsize\bfseries\color{darkblue}}
{\thesubsection}{0.75em}{}
\titleformat{\subsubsection}
{\normalsize\bfseries\color{medblue}}
{\thesubsubsection}{0.75em}{}

\pagestyle{fancy}
\fancyhf{}
\fancyhead[L]{\small\color{midgrey}ICV 442 Acueductos y Alcantarillados}
\fancyhead[R]{\small\color{midgrey}\DocShort}
\fancyfoot[C]{\small\color{midgrey}\thepage}
\renewcommand{\headrulewidth}{0.4pt}
\renewcommand{\footrulewidth}{0pt}

\rowcolors{2}{lightgrey}{white}
\renewcommand{\arraystretch}{1.18}

\newtcolorbox{infobox}{
  breakable,
  colback=lightblue,
  colframe=medblue,
  boxrule=0.8pt,
  arc=3pt,
  left=8pt,
  right=8pt,
  top=7pt,
  bottom=7pt
}

\newtcolorbox{warnbox}{
  breakable,
  colback=warnyellow,
  colframe=orange!70!black,
  boxrule=0.8pt,
  arc=3pt,
  left=8pt,
  right=8pt,
  top=7pt,
  bottom=7pt
}

\newcommand{\doccover}[3]{
  \begin{center}
  {\small\color{midgrey}Pontificia Universidad Cat\'olica Madre y Maestra\\Escuela de Ingenier\'ia Civil y Ambiental}

  \vspace{1.0cm}
  {\color{medblue}\rule{\linewidth}{2pt}}
  \vspace{0.3cm}

  {\Large\bfseries\color{darkblue}ICV-442-T Acueductos y Alcantarillado}\\[0.35em]
  {\LARGE\bfseries\color{darkblue}#1}\\[0.25em]
  {\large\itshape #2}

  \vspace{0.3cm}
  {\color{medblue}\rule{\linewidth}{2pt}}
  \vspace{0.9cm}
  \end{center}

  \begin{center}
  \begin{tabular}{ll}
  \toprule
  \textbf{Asignatura} & ICV-442-T Acueductos y Alcantarillado \\
  \textbf{Profesor} & Ricardo Rafael Hern\'andez Moreira, PhD \\
  \textbf{Periodo acad\'emico} & Mayo--agosto de 2026 \\
  \textbf{Proyecto} & Dise\~no de sistemas de saneamiento para Las Charcas, Azua \\
  \textbf{Versi\'on} & #3 \\
  \bottomrule
  \end{tabular}
  \end{center}

  \vspace{0.4cm}
}


\begin{document}
\doccover{R\'ubrica de evaluaci\'on}{Proyecto escrito y presentaci\'on oral}{18 de mayo de 2026}

\section{Proyecto escrito}
\small
\begin{longtable}{p{0.20\linewidth} p{0.16\linewidth} p{0.16\linewidth} p{0.16\linewidth} p{0.16\linewidth} p{0.07\linewidth}}
\toprule
\textbf{Criterio} & \textbf{Excelente} & \textbf{Satisfactorio} & \textbf{En desarrollo} & \textbf{Insuficiente} & \textbf{Peso} \\
\midrule
\endfirsthead
\toprule
\textbf{Criterio} & \textbf{Excelente} & \textbf{Satisfactorio} & \textbf{En desarrollo} & \textbf{Insuficiente} & \textbf{Peso} \\
\midrule
\endhead
Calidad de la informaci\'on b\'asica & Fuentes oficiales claras, proyecci\'on bien justificada y c\'alculos trazables. & Informaci\'on suficiente con algunas debilidades de trazabilidad o justificaci\'on. & Informaci\'on parcial o poco sustentada. & Fuentes d\'ebiles o ausencia de base suficiente. & 15\% \\
Fundamentos hidr\'aulicos & Supuestos, c\'alculos y modelos coherentes; resultados verificados. & Trabajo correcto con omisiones menores. & Inconsistencias parciales o verificaci\'on incompleta. & Errores graves de criterio o falta de verificaci\'on. & 35\% \\
Cumplimiento normativo y justificaci\'on t\'ecnica & Par\'ametros clave bien sustentados y decisiones defendibles. & Mayormente sustentado, con algunos vac\'ios. & Justificaci\'on superficial o dispersa. & Criterios adoptados sin sustento suficiente. & 25\% \\
Claridad y presentaci\'on & Documento ordenado, legible y bien estructurado. & Presentaci\'on adecuada con detalles por mejorar. & Estructura confusa en varias partes. & Documento desordenado o dif\'icil de seguir. & 15\% \\
Viabilidad y especificaciones & Materiales, detalles y mantenimiento consistentes con el contexto. & En general viable, aunque incompleto en algunos puntos. & Viabilidad poco discutida o mal cerrada. & Sin consideraci\'on seria de viabilidad o mantenimiento. & 10\% \\
\bottomrule
\end{longtable}
\normalsize

\section{Presentaci\'on oral}
\begin{center}
\begin{tabular}{p{0.34\linewidth} p{0.49\linewidth} p{0.08\linewidth}}
\toprule
\textbf{Criterio} & \textbf{Se espera} & \textbf{Peso} \\
\midrule
S\'intesis y claridad & Exposici\'on ordenada, dentro del tiempo, con lenguaje t\'ecnico claro. & 30\% \\
Dominio t\'ecnico & Capacidad para explicar criterios, defender supuestos y responder preguntas. & 40\% \\
Material visual & Gr\'aficos, planos y resultados legibles, sin sobrecarga. & 15\% \\
Trabajo en equipo & Participaci\'on equilibrada y transiciones razonables. & 15\% \\
\bottomrule
\end{tabular}
\end{center}

\section{Observaciones de uso}
\begin{itemize}
\item La r\'ubrica se aplica tanto al producto final como al nivel de control t\'ecnico demostrado por el equipo.
\item Un documento muy extenso no compensa la falta de criterio en decisiones clave.
\item La calidad de la presentaci\'on oral no sustituye debilidades del proyecto escrito, ni a la inversa.
\end{itemize}

\end{document}
